\section{Related Work}
Four strands of earlier work are relevant here: UAV-assisted vehicular edge
computing, queue-aware stochastic control, multi-agent reinforcement learning
and robust policy training. We review each one with an eye on what it offers
for task service and control in our setting.

\subsection{UAV-Assisted MEC and End-to-End Task Service}
Putting an edge server on a UAV shortens the access distance and lets the
server follow demand as it moves around, but it also ties the trajectory,
radio service, computation and energy budget to one another. Joint
formulations of offloading, CPU control, wireless power delivery and UAV
motion deal with this coupling directly~\cite{liu2020uav,xiong2021energy}.
With several UAVs, cooperative offloading and resource allocation must also
respect shared airspace and capacity limits~\cite{guo2024multi}, and
multi-agent policies have recently been used for task partitioning and
resource allocation in vehicular MEC~\cite{liang2025rav}. These papers make a
good case for joint control. Even so, they often describe a task by a single
delay term or by an allocation that lasts one slot.

Our model follows a task through the stages it actually visits (ingress,
radio upload, receive staging, edge computation, UAV--cloud backhaul, cloud
computation and branch synchronization), so waiting that happens after the
radio link is counted. A task that misses its deadline is kept in the
service model until it finishes or the trace ends.

\subsection{Queue-Aware Vehicular Offloading and Lyapunov Control}
Lyapunov optimization reduces a long-horizon stochastic control problem to
queue-weighted decisions taken slot by slot, provided that arrivals, service
and energy increments meet the boundedness and feasibility conditions it
needs~\cite{neely2010stochastic}. Mao \emph{et al.} applied it to MEC
offloading with energy-harvesting devices~\cite{mao2017dynamic}. For
vehicular edge computing, Q-TOP couples physical queues with 3-D path
planning for multiple UAVs and solves each per-slot problem with a genetic
optimizer~\cite{sun2026qtop}.

We keep the queue-weighted structure and extend it with explicit transport,
staging, backhaul and cloud queues, per-task identities and a virtual
net-energy deficit. Moving work on to the next stage does not count as
finishing it, and the drift term registers congestion both at the queue the
work leaves and at the one it enters.

\subsection{Multi-Agent and Hybrid-Action Reinforcement Learning}
MADDPG trains centralized critics but executes decentralized actors, which
lets agents with different observations pursue one shared
objective~\cite{lowe2017multi}. MAPPO is an on-policy alternative that keeps
each policy update small through clipping~\cite{schulman2017proximal}. An
attention-based MAPPO has been used to coordinate bounded UAV motion, task
service and energy transfer~\cite{talukdar2026mappo}, and generative policies
have been tried for UAV task assignment~\cite{tang2025dnn}.

Our action space mixes discrete and continuous parts. A vehicle decides
whether to hold a task or where to send it, and it also picks a continuous
partition fraction. RSUs and UAVs set continuous processor shares, and UAVs
additionally choose a CPU frequency and a 3-D motion. The controller copes
with this through masked straight-through categorical outputs, bounded
continuous heads and a deterministic feasibility layer. Neither MADDPG nor
PPO gives battery causality, queue stability or collision avoidance on its
own.

\subsection{Meta-Learning and Robustness Across Operating Domains}
MAML and Reptile learn an initialization that should adapt after a few
gradient steps~\cite{finn2017model,nichol2018reptile}, and Yan \emph{et al.}
combine Reptile with MADDPG for computation allocation in a UAV/RSU/cloud MEC
system~\cite{yan2026rmaddpg}. Domain randomization has another aim. A single
policy is optimized across a distribution of environments and then deployed
with fixed parameters, an approach that has made learned controllers more
robust to model variation~\cite{tobin2017domain,peng2018sim}.

This affects how our results should be read. Training covers a declared
family of load, harvesting, outage and backhaul conditions, and the frozen
policy is then tested on separate stochastic traces from the same family. The
experiments therefore measure how a fixed policy generalizes. They do not
demonstrate few-shot adaptation, nor robustness to operating changes outside
that family.

\subsection{Energy-Aware Flight and Safe Execution}
Propulsion, communication and computation all draw on one finite battery, so
UAV-MEC studies usually model rotary-wing flight together with the offloading
and trajectory decisions~\cite{liu2020uav,talukdar2026mappo}. Work on
predictive UAV deployment likewise links position to both demand and
mobility energy~\cite{zhang2019predictive}. In our design, immediate battery
causality and a verified return/landing reserve are hard execution
constraints, and a virtual deficit queue tracks the declared long-run
net-energy target. Collisions and battery exhaustion are prevented by having
the safety controller check whole slot trajectories; we do not rely on a
reward penalty for this.

\subsection{Synthesis and Research Gap}
Table~\ref{tab:positioning} sums up how these strands relate to our design.
Mobile edge servers, Lyapunov queue control, multi-agent learning, robust
training and energy-aware flight have each proved useful, but they do not
simply stack on top of each other. Routing changes what arrives at a
downstream queue. A processor or backhaul allocation changes delay and energy
at once. UAV motion changes the future channel, the propulsion demand and the
return reserve. If one of these effects is modeled outside the shared state
and transition, congestion or energy cost can be pushed onto a stage that the
controller never sees.

We see three gaps. First, a large part of the offloading literature describes
service by an instantaneous rate or an aggregate delay. That hides where
dispatched but unfinished work is waiting, and an evaluation built on it can
count work as done while its upload, computation, synchronization or backhaul
is still pending. Second, MADDPG and MAPPO provide a learning rule but do not
enforce processor capacity, battery causality, continuous-path separation or
a safe return reserve, and reward penalties cannot guarantee these at
execution time. Third, a policy tuned in a nominal simulator may not stay
reliable when load, harvesting, outages or backhaul change. Meta-learning
targets fast adaptation, but it assumes that data and gradient updates are
available after deployment, and our zero-update setting rules that out.

What is needed is a fixed multi-agent policy that observes end-to-end
congestion, coordinates heterogeneous agents with mixed actions, respects hard
physical limits, and is trained and tested on a declared family of operating
regimes. We address this with persistent task-level queues, physical and
virtual queue signals derived from Lyapunov analysis, hybrid CTDE control and
a deterministic feasibility layer. During evaluation the trained actors are
frozen and run on paired held-out traces, so differences between controllers
cannot come from different traffic or channel realizations. Our claim is
accordingly limited to robustness within the declared domain family; we do
not claim guaranteed performance under arbitrary out-of-distribution
conditions.

\begin{table*}[!t]
\caption{Positioning of the Proposed Framework Relative to the Main Research Streams}
\label{tab:positioning}
\centering
\footnotesize
\begin{tabular}{p{0.19\textwidth}p{0.34\textwidth}p{0.39\textwidth}}
\toprule
Research stream & Established capability & Gap addressed in this work\\
\midrule
UAV-assisted MEC~\cite{liu2020uav,xiong2021energy,guo2024multi} & Joint optimization of offloading, computing, radio allocation, and trajectory-aware service & Stage queues along the vehicle--RSU--UAV--cloud path, per-task records, branch synchronization, and a finite backhaul\\
Queue-aware VEC~\cite{neely2010stochastic,mao2017dynamic,sun2026qtop} & Offloading and UAV path planning weighted by Lyapunov queues & Transport and staging queues, weighted CPU sharing, censoring of unfinished tasks, and a virtual net-energy queue\\
MADDPG/MAPPO~\cite{lowe2017multi,schulman2017proximal,talukdar2026mappo} & Coordinated multi-agent policies, deterministic or stochastic & Masked categorical routing together with continuous split, CPU, and flight actions\\
Meta-RL~\cite{finn2017model,nichol2018reptile,yan2026rmaddpg} & An initialization that adapts with few updates & Deployment with fixed parameters and zero updates; the evaluated policy makes no few-shot claim\\
Domain randomization~\cite{tobin2017domain,peng2018sim} & Robust control over a declared distribution of parameters & Paired held-out tests across traffic, harvesting, outage, and backhaul regimes\\
Energy-aware UAV control~\cite{liu2020uav,talukdar2026mappo,zhang2019predictive} & Models of propulsion and computation energy & Energy ledger for the whole system, battery causality, collision avoidance, and a verified return/landing reserve\\
\bottomrule
\end{tabular}
\end{table*}
